\section*{Chapitre \ref{ch:iv1:timed-test-9} --- Épreuve chronométrée 9}
\markboth{\protect\omTocMark{}{Réponses~: Épreuve chronométrée 9}}{\protect\omTocMark{}{Réponses~: Épreuve chronométrée 9}}

\begin{solution}{iq:iv1:timed-test-9:1}
L'espérance cherchée vaut $\frac83$.

Soient $G_1$ et $G_2$ les rangs du premier pile de chaque pièce. On veut $\mathbb{E}\max(G_1, G_2)$. En utilisant $\max + \min = G_1 + G_2$\index{maximum et minimum}, la somme a pour espérance 4.

Le minimum est le premier tour où l'une des deux pièces montre pile. Un tour réalise cet événement avec la probabilité $\frac34$, donc le minimum a pour espérance $\frac43$. Le maximum a donc pour espérance $4 - \frac43 = \frac83$.

Une seconde voie passe par les états~: depuis le départ, $E = 1 + \frac14E + \frac12 \cdot 2$, ce qui donne $\frac34E = 2$, donc $E = \frac83$. Une simulation confirme ce résultat.

\iqwrong{Répondre 4 en additionnant les deux espérances, ou 2 en supposant qu'elles avancent en parallèle. La vraie valeur se situe entre les deux.}
\end{solution}

\begin{solution}{iq:iv1:timed-test-9:2}
La probabilité vaut $\frac17$.

Il y a $\binom83 = 56$ triplets. Tout triangle équilatéral doit avoir trois côtés égaux, et les distances entre sommets sont l'arête, la diagonale de face et la diagonale de l'espace.

Seule la diagonale de face convient\index{comptage par symétrie}. Chaque sommet a 3 autres sommets à distance de diagonale de face, et ces trois sommets sont aussi deux à deux à cette distance. Donc deux quelconques d'entre eux forment un triangle équilatéral avec le sommet, et chaque sommet appartient à $\binom32 = 3$ triangles. Chaque triangle ayant 3 sommets, il y a $\frac{8 \times 3}{3} = 8$ triangles. Alors $\frac{8}{56} = \frac17$.

Une énumération des 56 triplets confirme le 8. Aucun triplet d'arêtes ne se referme, car un cube ne contient pas de triangle formé par ses arêtes.

\iqwrong{Compter seulement 2 triangles, ou les triangles d'un tétraèdre, dont il y a 4. Le cube en contient deux de ce type.}
\end{solution}

\begin{solution}{iq:iv1:timed-test-9:3}
La moyenne a posteriori vaut 1 et la variance $\frac12$. Avec une variance de bruit égale à 4, elles valent $0.4$ et $0.8$.

Les précisions s'additionnent\index{a priori normal conjugué}. L'a priori a une précision de 1, et la mesure a une précision de 1. La précision a posteriori vaut donc 2, la variance $\frac12$, et la moyenne est la moyenne pondérée par les précisions~: $\frac{1 \cdot 0 + 1 \cdot 2}{2} = 1$.

Avec une variance de bruit égale à 4, la mesure a une précision de $\frac14$. La précision a posteriori devient $\frac54$, la variance vaut $0.8$, et la moyenne $\frac{\frac14 \cdot 2}{\frac54} = 0.4$.

Une intégration numérique de la loi a posteriori confirme ces valeurs. Vérifications aux limites~: sans bruit, la moyenne vaut 2, et avec un bruit infini, elle reste à 0.

\iqwrong{Donner une variance a posteriori supérieure à celle de l'a priori. Dans ce cas normal, une nouvelle mesure ne peut que la réduire.}
\end{solution}

\begin{solution}{iq:iv1:timed-test-9:4}
La probabilité vaut $\frac{16}{49}$.

Comme $10^6 = 2^6 5^6$, les diviseurs sont les $2^a5^b$ avec $a$ et $b$ chacun entre 0 et 6, ce qui donne $7 \times 7 = 49$\index{comptage des diviseurs}.

Un multiple de $10^3$ exige $a \ge 3$ et $b \ge 3$. Cela fait 4 valeurs pour chaque exposant, soit 16 diviseurs. D'où $\frac{16}{49}$.

Une liste directe des diviseurs confirme les décomptes. La probabilité de $a \ge 3$ seul vaut $\frac47$, et les deux exposants sont indépendants, donc la réponse est son carré. Notons le couplage~: un diviseur $d$ correspond à $10^6/d$, et $10^3$ se trouve au milieu de la liste.

\iqwrong{Utiliser la probabilité $\frac{1}{1000}$, comme pour un entier pris au hasard. Les diviseurs ne sont pas répartis uniformément, donc ce sont les exposants qu'il faut compter.}
\end{solution}

\begin{solution}{iq:iv1:timed-test-9:5}
La probabilité vaut environ $0.24$, et l'espérance conditionnelle vaut 1.

$X - Y$ suit une loi normale d'espérance 0 et de variance 2\index{somme de lois normales}. Donc $P(X - Y > 1) = P(Z > \frac{1}{\sqrt2}) = P(Z > 0.71) = 0.24$.

Sachant $X + Y = 2$, la symétrie donne $\mathbb{E}[X \mid X + Y = 2] = \mathbb{E}[Y \mid X + Y = 2]$, et les deux s'additionnent à 2. Chacune vaut donc 1.

La voie par la régression confirme ce résultat~: la pente de $X$ sur $X + Y$ vaut $\frac12$, et la moyenne de $X + Y$ est 0, donc la valeur vaut $\frac12 \times 2$. Une simulation donne les deux résultats.

Vérification aux limites~: pour une cible de 0 au lieu de 1, la probabilité vaut $\frac12$, comme il se doit.

\iqwrong{Prendre une variance de 1 pour $X - Y$. Les variances s'additionnent aussi pour une différence, donc l'échelle est $\sqrt2$.}
\end{solution}

\begin{solution}{iq:iv1:timed-test-9:6}
Le plus petit a pour espérance $\frac12$, le plus grand pour espérance $\frac32$, et la probabilité vaut $\frac{1}{2e} \approx 0.184$.

Le minimum suit une loi exponentielle de paramètre 2\index{absence de mémoire}, donc son espérance vaut $\frac12$. Après lui, l'autre variable repart d'une durée de vie exponentielle neuve, par absence de mémoire. L'écart a donc pour espérance 1, et le maximum a pour espérance $\frac32$.

La somme se vérifie~: $\frac12 + \frac32 = 2 = \mathbb{E}(X + Y)$.

$X > Y + 1$ suppose deux conditions~: $X$ survit à $Y$, avec la probabilité $\frac12$, puis il dure au moins 1 de plus, avec la probabilité $e^{-1}$. Une simulation donne $0.184$.

Vérification aux limites~: remplacer 1 par 0 donne $\frac12$, et un grand écart annule la probabilité.

\iqwrong{Supposer que l'écart dépend de la taille de $Y$. L'absence de mémoire le rend indépendant de la durée de vie déjà écoulée du premier.}
\end{solution}

\begin{solution}{iq:iv1:timed-test-9:7}
$P(N \text{ pair}) = \frac{1 + e^{-2}}{2} \approx 0.568$, et la probabilité de ne faire aucune vente vaut $e^{-3} \approx 0.050$.

Additionner les séries de $e^{\lambda}$ et de $e^{-\lambda}$\index{astuce de la parité}~: les termes impairs s'annulent. Donc $\sum_{k \text{ pair}}\frac{\lambda^k}{k!} = \cosh\lambda$, et $P(\text{pair}) = e^{-\lambda}\cosh\lambda = \frac{1 + e^{-2\lambda}}{2}$. Pour $\lambda = 1$, cela vaut $0.568$. Quand $\lambda$ croît, la probabilité tend vers $\frac12$, et pour $\lambda = 0$ elle vaut 1.

L'amincissement d'un processus de Poisson\index{amincissement de Poisson} avec une probabilité de $\frac12$ donne un processus de Poisson de paramètre 3. La probabilité de ne faire aucune vente vaut donc $e^{-3}$. Sommer sur le nombre de clients donne le même résultat~: $\sum e^{-6}\frac{6^k}{k!}2^{-k} = e^{-6}e^{3}$. Une simulation confirme ce résultat.

\iqwrong{Utiliser $e^{-6}$, qui est la probabilité de n'avoir aucun client. Un client qui n'achète pas reste une absence de vente.}
\end{solution}

\begin{solution}{iq:iv1:timed-test-9:8}
La probabilité vaut $\frac{1800}{7776} = \frac{25}{108} \approx 0.231$.

Choisir les deux valeurs formant la paire~: $\binom62 = 15$. Choisir la valeur isolée~: 4. Ordonner les cinq dés\index{coefficient multinomial}~: $\frac{5!}{2!\,2!\,1!} = 30$. Cela fait $15 \times 4 \times 30 = 1800$, sur $6^5 = 7776$ issues possibles.

Une énumération des 7776 issues confirme le décompte.

Le choix des deux valeurs de la paire se fait sans ordre. Le faire avec ordre compterait chaque issue deux fois.

\iqwrong{Utiliser $6 \times 5 \times 4 \times 30 = 3600$, ce qui traite (3, 3, 5, 5, 2) et (5, 5, 3, 3, 2) comme deux résultats différents.}

\iqfollow{Est-ce plus ou moins probable qu'une seule paire~?}{Une seule paire compte $6 \times \binom53 \times \frac{5!}{2!} = 3600$ issues, soit $0.463$, donc elle est deux fois plus probable.}
\end{solution}

\begin{solution}{iq:iv1:timed-test-9:9}
Les poids valent 2 pour 1, soit $\frac23$ et $\frac13$, et le ratio de Sharpe vaut environ $0.707$.

Le meilleur portefeuille est proportionnel à $\Sigma^{-1}\mu$\index{portefeuille de Sharpe maximal}. Sans corrélation, les poids sont proportionnels à $\mu_i/\sigma_i^2$~: $\frac{0.05}{0.01} = 5$ et $\frac{0.10}{0.04} = 2.5$. Cela fait 2 pour 1.

Chaque actif seul a un ratio de Sharpe de $\frac12$. Pour des actifs non corrélés, les carrés des ratios de Sharpe s'additionnent, donc le meilleur vaut $\sqrt{\frac14 + \frac14} = 0.707$.

Vérification directe~: le portefeuille a un rendement de $\frac23 \cdot 5 + \frac13 \cdot 10 = 6.67~\%$, et une variance de $\frac49 \cdot 100 + \frac19 \cdot 400 = 88.9$, donc la volatilité vaut $9.43~\%$. Le ratio vaut $\frac{6.67}{9.43} = 0.707$.

\iqwrong{Choisir des poids égaux, puisque les deux ratios de Sharpe sont identiques. Des ratios de Sharpe égaux supposent des risques égaux, et non des capitaux égaux.}
\end{solution}

\begin{solution}{iq:iv1:timed-test-9:10}
La probabilité vaut environ $0.0475$.

Choisir les deux valeurs~: $\binom{13}{2} = 78$. Choisir les deux couleurs dans chaque main\index{comptage des mains}~: $\binom42^2 = 36$. Choisir la cinquième carte parmi les 44 cartes des autres valeurs. Le compte vaut $78 \times 36 \times 44 = 123\,552$. Divisé par le nombre de mains, cela donne $0.0475$.

Une simulation donne le même résultat. Remarquer le 44~: seules les cartes des 11 autres valeurs sont permises, et chacune existe en 4 couleurs.

Pour une vérification rapide, $78 \times 36 \times 44 \approx 80 \times 1600 = 128\,000$, et ce nombre rapporté à $2.6$ million vaut environ une main sur 20.

\iqwrong{Utiliser $13 \times 12$ pour les valeurs, ce qui compte chaque main deux fois. Les deux paires n'ont pas d'ordre, donc la valeur double pour atteindre environ $0.095$.}
\end{solution}

\begin{solution}{iq:iv1:timed-test-9:11}
Le dernier passager obtient sa propre place avec probabilité $\frac12$, et le neuvième est déplacé avec probabilité $\frac13$.

À chaque fois qu'une personne s'assoit au hasard, trois issues sont possibles\index{symétrie}. Elle peut choisir la place du premier passager, et tout est réglé. Elle peut choisir la place du dernier, et le dernier est alors déplacé. Ou bien elle choisit une autre place, et le choix passe à quelqu'un d'autre. La première et la deuxième issues sont équiprobables à chaque étape, donc la chance du dernier passager vaut $\frac12$.

Pour le neuvième, le même raisonnement vaut avec trois places encore en jeu~: la place 1, la place 9 et la place 10. Lorsqu'un choix aléatoire tombe d'abord sur l'une d'elles, chacune est équiprobable. Seule la place 9 déplace le neuvième passager. La chance vaut donc $\frac13$. En général, le passager $k$ parmi $n$ est déplacé avec une chance de $\frac{1}{n - k + 2}$.

Une simulation confirme les deux résultats. Avec $k = 10$, on obtient $\frac12$, comme précédemment.

\iqwrong{Calculer l'arbre complet des 10 passagers. La réponse pour le dernier passager ne dépend pas de $n$ du tout, et c'est l'indice qu'il faut chercher une symétrie.}
\end{solution}

\begin{solution}{iq:iv1:timed-test-9:12}
$100!$ a 158 chiffres, et la limite vaut $\frac1e \approx 0.368$.

Avec $n = 100$, $n\ln n = 100 \times 4.605 = 460.5$\index{formule de Stirling}. En retranchant 100, on obtient $360.5$. Et $\frac12\ln(628) = 3.22$, donc $\ln 100! \approx 363.7$. En divisant par $2.303$, on obtient environ $157.9$. Sa partie entière vaut 157, donc $100!$ a 158 chiffres. Un produit exact le confirme~: $100! \approx 9.3 \times 10^{157}$.

La réponse est proche de la puissance de dix suivante, donc une estimation soignée compte ici~: le terme $\frac12\ln(2\pi n)$ est indispensable.

Pour la limite, $\frac1n\ln n! \approx \ln n - 1$. Donc $(n!)^{1/n} \approx \frac{n}{e}$. Cela signifie que la moyenne géométrique de $1, \dots, n$ vaut environ $\frac{n}{e}$, contre une moyenne arithmétique d'environ $\frac{n}{2}$.

\iqwrong{Négliger le petit terme et obtenir $\frac{360.5}{2.303} = 156.5$, soit 157 chiffres. C'est trop bas d'une unité.}
\end{solution}
